\section*{第一卷内容索引}

\begin{supertabular}{@{}>{\raggedright\arraybackslash}p{0.84\linewidth}>{\raggedleft\arraybackslash}p{0.10\linewidth}@{}}
\multicolumn{2}{c}{\textsc{第一章. 双曲率曲线。}}\\[0.3em]
§. 1. 切线与法平面 & 1\\
2. 第一曲率或弯曲率 & 2\\
3. 密切平面 & 4\\
4. 主三面形 & 6\\
5. 第二曲率或挠率 & 9\\
6. Frenet 公式 & 10\\
7. 挠率的符号 & 12\\
8--9. 内在方程 & 13\\
10--11. 柱面螺线 & 18\\
12. 包络面 & 22\\
13. 管状面 & 25\\
14. 可展曲面 & 26\\
15. 矫直可展面与可展曲面的测地线 & 27\\
16. 极可展面 & 31\\
17. 密切球球心的轨迹 & 32\\
18. 渐开线与渐屈线 & 35\\
19. 一个 \(\infty^1\) 平面族的正交轨迹 & 38\\
20. 曲线的 Combescure 变换 & 40\\
21. 第一曲率恒定的曲线与球面螺线 & 42\\
22--23. 一个 \(\infty^1\) 球面族的正交轨迹 & 44\\
24--25. Bertrand 曲线 & 50\\[0.5em]
\multicolumn{2}{c}{\textsc{第二章. 二次微分形式。}}\\[0.3em]
§. 26. 代数二次形式 & 54\\
27. 二次微分形式 & 58\\
28. 不变量与微分参数 & 59\\
29. 第一微分参数 \(\Delta_1U,\Delta(U,V)\) & 60\\
30. 二次型的等价性与 Christoffel 公式 & 62\\
31. 三指标符号及其性质 & 64\\
32. 第二微分参数 \(\Delta_2U\) & 66\\
33. 对 \(P_{ik}\) 的全微分方程 & 69\\
34. 可积性条件 & 71\\
35. 四指标 Riemann 符号及其性质 & 72\\
36. 无限可积的情形 & 74\\
37. 二元型的曲率 & 76\\
38--39. 联立二次型 & 79\\[0.5em]
\multicolumn{2}{c}{\textsc{第三章. 曲面上的曲线坐标。——共形表示。}}\\[0.3em]
§. 40. 曲线坐标 & 84\\
41. 线元和面积元 & 86\\
42. 曲线同坐标线的夹角 & 89\\
43. Christoffel 符号、微分参数与曲率的值 & 92\\
44. 等温系 & 94\\
45. 零长度线 & 96\\
46. 等距参数 & 98\\
47. Lie 关于等温系统的定理 & 100\\
48--49. 共形表示 & 102\\
50. 对旋转曲面的应用 & 105\\
51. 复球面 & 106\\
52. 平面或球面上的正交双重圆系 & 108\\
53. 复球面的运动（Cayley 公式） & 109\\[0.5em]
\multicolumn{2}{c}{\textsc{第四章. 曲面理论的基本公式。}}\\[0.3em]
§. 54. 曲面的两个基本二次型 & 112\\
55. 基本方程 & 115\\
56. Gauss 方程与 Codazzi 方程 & 117\\
57--58. 曲面的内在方程及其积分 & 121\\
59. 曲率线 & 126\\
60. 曲率线的微分方程 & 127\\
61. 法截线的曲率——Meusnier 定理 & 129\\
62. Euler 公式 & 131\\
63. 平均曲率与总曲率 & 133\\
64. 共轭切线 & 135\\
65. 渐近线 & 138\\
66--67. 共轭系的性质与例子 & 139\\
68. 笛卡儿坐标中的曲率线与主曲率半径 & 144\\
69. 微分参数计算 & 145\\[0.5em]
\multicolumn{2}{c}{\textsc{第五章. Gauss 球面表示。——切向坐标。}}\\[0.3em]
§. 70. Gauss 球面表示 & 148\\
71. 渐近线挠率与 Enneper 定理 & 150\\
72. 球面表示一般公式 & 152\\
73. 以渐近线为坐标的曲面 & 155\\
74. 渐近线上挠率符号 & 157\\
75. 面积极小曲面上的渐近线 & 159\\
76. 伪球曲面上的渐近线 & 160\\
77. 渐近线的 Lelieuvre 公式 & 162\\
78. 指定共轭系球面像的曲面 & 166\\
79. 等温共轭系 & 167\\
80. 新的 Lelieuvre 公式 & 170\\
81. 切向坐标中的 Weingarten 公式 & 172\\
82. 切向坐标中共轭系的性质 & 174\\
83. 具有位于平行平面内的曲率线系的曲面 & 176\\[0.5em]
\multicolumn{2}{c}{\textsc{第六章. 测地曲率。——测地线。}}\\[0.3em]
§. 84. 曲线的切向曲率 & 179\\
85. Bonnet 关于切向曲率（测地曲率）的公式 & 181\\
86. 曲率 K 的 Liouville 形式 & 183\\
87. 测地线及其微分方程 & 185\\
88. Gauss 关于测地线的微分方程 & 189\\
89. 测地线的一般性质 & 191\\
90. 线元的测地形式 & 194\\
91. 测地椭圆与测地双曲线 & 197\\
92. 测地挠率 & 198\\
93. 测地线方程积分的一般定理 & 202\\
94. Liouville 曲面上的测地线 & 205\\
95. 旋转曲面的测地线 & 207\\
96. 测地三角形的总曲率（Gauss 定理） & 208\\
97. 常测地曲率曲线的正交双重系统 & 210\\[0.5em]
\multicolumn{2}{c}{\textsc{第七章. 等距适用曲面。}}\\[0.3em]
§. 98. 可弯曲而不可伸长的曲面——等距适用曲面 & 213\\
99. 第一个等距适用问题及其解 & 217\\
100. 常曲率曲面的情形 & 219\\
101. 常曲率曲面在自身上的 \(\pxIndexConstantCurvatureSelfApplicability\) 个等距适用 & 222\\
102. 伪球面 \(ds^2\) 的三种标准形式 & 223\\
103. 旋转伪球面 & 225\\
104. 在自身上有 \(\infty^1\) 个等距适用的曲面 & 230\\
105. 彼此等距适用的旋转曲面 & 231\\
106. 螺旋面与 Bour 定理 & 235\\
107. 面积极小直纹螺旋面与单叶双曲面 & 236\\
108--109. 第二个等距适用问题及其所依赖的偏微分方程 & 238\\
110. 包含一条刚性曲线的曲面的不可变形性 & 241\\
111--112. 使一条给定曲线取得指定形状的变形 & 244\\
113--114. 特殊变形 & 249\\[0.5em]
\multicolumn{2}{c}{\textsc{第八章. 直纹曲面的变形。}}\\[0.3em]
§. 115. 关于等距适用直纹曲面的 Bonnet 定理 & 252\\
116. 直纹曲面的线元 & 254\\
117. 腰线 & 256\\
118. 切平面分布参数 & 257\\
119. 按照 Minding 方法变形直纹曲面 & 260\\
120. Beltrami 方法 & 262\\
121--122. 直纹曲面的特殊变形 & 264\\
123. 与旋转曲面等距适用的直纹曲面 & 268\\[0.5em]
\multicolumn{2}{c}{\textsc{第九章. 演化曲面与 Weingarten 定理。}}\\[0.3em]
§. 124. 曲面演化曲面的两个分支 & 271\\
125--126. 演化曲面基本量的计算 & 273\\
127. Beltrami 关于测地曲率半径的构造 & 277\\
128. Ribaucour 的平均演化曲面 & 278\\
129. W 曲面与演化曲面的两个分支 & 280\\
130. Ribaucour 关于曲率线对应关系的定理 & 282\\
131. Lie 关于 W 曲面的定理 & 284\\
132. Weingarten 关于 W 曲面演化曲面的第一定理 & 285\\
133. Weingarten 的第二定理 & 287\\
134. W 曲面的代表球面线元 & 289\\
135. W 曲面的完备类别 & 291\\
136. 伪球曲面的演变曲面与互补曲面 & 293\\
136*. 互补曲面的线元计算 & 295\\[0.5em]
\multicolumn{2}{c}{\textsc{第十章. \(\infty^2\) 条射线系或直线汇。}}\\[0.3em]
§. 137. 直线汇的两个基本二次型 & 297\\
138. 极限点与 Hamilton 公式 & 300\\
139. Ribaucour 各向同性直线汇 & 302\\
140. 极限点的横坐标 & 304\\
141. 直线汇的焦点与可展曲面 & 305\\
142. 焦曲面 & 307\\
143. 法向直线汇 & 309\\
144. Malus--Dupin 定理 & 310\\
145. 焦曲面具有指定球面像的直线汇 & 312\\
146. 各向同性直线汇的平均包络面 & 314\\
147. 可展曲面具有指定球面像的直线汇 & 316\\
148. 球面相关的一般公式 & 318\\
149. 焦曲面两个分支的元素 & 321\\
150. 对伪球直线汇的应用 & 323\\
151. Guichard 直线汇 & 325\\[0.5em]
\multicolumn{2}{c}{\textsc{第十一章. \(n\) 维空间。——这些空间的超曲面与曲线。}}\\[0.3em]
§. 152. \(n\) 维空间，以其 \(ds^2\) 定义 & 327\\
153. 角度度量 & 329\\
154. 方向常数 & 331\\
155. 测地线 & 333\\
156. \(ds^2\) 的测地形式 & 336\\
157. 任意方向上的 Riemann 曲率 & 339\\
158. Riemann 曲率公式 & 340\\
159. 常曲率 \(n\) 维空间 & 343\\
160. 具有同一常曲率的两个空间之间的等距适用性 & 346\\
161. Schur 定理 & 349\\
162. 三维空间中的三个主曲率 & 352\\
163. \(S_3\) 中的主三面体 & 355\\
164. 曲空间中的超曲面及其两个基本二次型 & 357\\
165. 推广的 Gauss 方程与 Codazzi 方程 & 359\\
166. 高维空间中的曲线 & 363\\
167. 超曲面上曲线的曲率 & 365\\
168. 超曲面的主曲率与推广的 Euler 公式 & 368\\
169. 特殊情形 & 371\\
170. 共形表示中的空间与 Liouville 定理 & 373\\
171. 正交 \(n\) 重系统在 \(n\) 维空间中的情形 & 377\\
172. 三维空间的情形——曲率线的性质与推广的 Dupin--Darboux 定理 & 379\\
173. \(S_3\) 中的三重正交系 & 383\\[0.5em]
\multicolumn{2}{c}{\textsc{第十二章. 二维伪球几何。}}\\[0.3em]
§. 174. 伪球面在上半平面上的共形表示 & 386\\
175. 等距适用变换或第一类运动 & 389\\
176. 第二类运动 & 391\\
177. 另一种共形表示——圆的保持 & 392\\
178. 平行测地线与平行角 & 395\\
179. 测地三角形与伪球三角学 & 397\\
180. 非欧几何与 Beltrami 表示 & 401\\
181. 由线性方程表示的测地线 & 403\\
182. 用于求测地线的 Riccati 方程 & 405\\
183. 非欧平面中的 Weierstrass 坐标 & 407\\
184. 关于 Weierstrass 坐标的 Weingarten 偏微分方程 & 410\\[0.5em]
\multicolumn{2}{c}{\textsc{第十三章. 常曲率空间几何。}}\\[0.3em]
§. 185. Weingarten 联立方程 & 413\\
186. 常曲率空间的 \(ds^2\) 典型形式 & 415\\
187. 伪球空间在欧几里得半空间上的共形表示 & 419\\
188. 测地线与测地子流形的像 & 421\\
189. 双曲 \(S_3\) 中的运动——Poincaré 公式 & 424\\
190. 双曲空间中的运动分类 & 426\\
191. Beltrami 测地表示 & 429\\
192. Cayley 度量 & 432\\
193. 椭圆度量与 Weierstrass 坐标 & 434\\
194. 双曲度量 & 437\\
195. 关于 Weierstrass 坐标的 Weingarten 方程 & 440\\
196. 椭圆几何中的滑移 & 444\\
197. 使绝对形变为自身的射影运动 & 447\\
198. Clifford 平行性 & 448\\
199. 滑移参数与平行角 & 450\\
200. 零曲率直纹曲面与 Clifford 曲面 & 453\\
201. 椭圆几何与双曲几何中的 Frenet 公式 & 455\\
202. 对常挠率曲线的应用 & 458\\[0.5em]
\multicolumn{2}{c}{\textsc{第十四章. 常曲率空间中的超曲面。}}\\[0.3em]
§. 203. 欧几里得空间中的超曲面 & 461\\
204. 欧几里得空间中超曲面的不可变形性 & 465\\
205. 等距适用性方程 & 468\\
206. 超曲面的曲率线与渐屈超曲面 & 471\\
207. 超曲面的球面表示 & 473\\
208. 切向坐标 & 478\\
209. 超球面中的正交 \(n\) 重系统与 Darboux 定理 & 481\\
210. 椭圆空间中超曲面的一般公式 & 482\\
211. 双曲空间中的超曲面公式 & 486\\
212. 椭圆几何与双曲几何中超曲面的曲率线 & 488\\
213. 三维椭圆几何与双曲几何中的曲面 & 491\\
214. 共轭系与渐近线 & 494\\
215. 以其曲率线为坐标的曲面 S & 497\\
216. 演化曲面的两个分支 & 499\\
217. W 曲面与 Weingarten 定理 & 501\\
218. 以渐近线为坐标的曲面 & 504\\
219. 椭圆几何中的零曲率曲面 & 506\\
220. 作为滑移曲面的零曲率曲面 & 508\\
221. 零曲率曲面的进一步性质 & 510\\
222. 曲空间中的曲面在欧几里得空间上作共形表示时的公式 & 513\\
\end{supertabular}
